\documentclass[aspectratio=169]{beamer}
\usetheme{metropolis}
\usepackage{appendixnumberbeamer}
\usepackage{booktabs, hyperref}

\input{mgmt675-style}

\usepackage{tikz}
\usetikzlibrary{shapes.geometric, arrows.meta, positioning, calc}

% Title info
\subtitle{MGMT 675: Generative AI for Finance}
\title{Module 2: Data Exploration \& Strategic Thinking}
\author{Kerry Back}
\date{}

\begin{document}

\maketitle

% ========================================
% SECTION 1: THE OAVR WORKFLOW
% ========================================
\section{The OAVR Workflow}

\begin{frame}{From Data to Decision: OAVR}
Most people jump straight to charts.  Better analysts follow a structured workflow: \alert{Orient $\rightarrow$ Ask $\rightarrow$ Visualize $\rightarrow$ Recommend}.
\vspace{0.3cm}
\begin{center}
\begin{tikzpicture}[
  node distance=1.2cm,
  every node/.style={font=\small},
  sbox/.style={draw=titlegray, rounded corners, minimum width=2.5cm, minimum height=0.9cm, fill=excelinput, text=titlegray, font=\footnotesize\bfseries, align=center},
  sarrow/.style={-{Stealth[length=3mm]}, thick, draw=accentblue}
]
\node[sbox] (orient) {Orient};
\node[sbox, right=of orient] (ask) {Ask};
\node[sbox, right=of ask] (viz) {Visualize};
\node[sbox, right=of viz] (rec) {Recommend};

\draw[sarrow] (orient) -- (ask);
\draw[sarrow] (ask) -- (viz);
\draw[sarrow] (viz) -- (rec);
\end{tikzpicture}
\end{center}
\vspace{0.3cm}
\begin{columns}[T]
\begin{column}{0.5\textwidth}
\begin{itemize}\small
  \item \textbf{Orient}: What decision are we making? Who is the audience? What data is available?
  \item \textbf{Ask}: What specific questions would change the decision?
\end{itemize}
\end{column}
\begin{column}{0.5\textwidth}
\begin{itemize}\small
  \item \textbf{Visualize}: Build charts and tables that answer those questions---nothing more
  \item \textbf{Recommend}: State a clear recommendation with supporting evidence
\end{itemize}
\end{column}
\end{columns}
\end{frame}

\begin{frame}{OAVR in Practice: Should We Expand to Austin?}
\begin{columns}[T]
\begin{column}{0.5\textwidth}
\begin{shadedbox}[title=\textbf{Orient}]
\begin{itemize}\scriptsize
  \item Decision: open a new office in Austin vs.\ Denver
  \item Audience: executive committee
  \item Data: real estate costs, talent pool, client proximity, tax incentives
\end{itemize}
\end{shadedbox}
\begin{shadedbox}[title=\textbf{Ask}]
\begin{itemize}\scriptsize
  \item What is the all-in cost per employee in each city?
  \item How many qualified hires are within 50 miles?
  \item What fraction of current clients are within a 2-hour flight?
\end{itemize}
\end{shadedbox}
\end{column}
\begin{column}{0.5\textwidth}
\begin{shadedbox}[title=\textbf{Visualize}]
\begin{itemize}\scriptsize
  \item Side-by-side cost comparison table
  \item Map of talent density by metro area
  \item Client accessibility heat map
\end{itemize}
\end{shadedbox}
\begin{shadedbox}[title=\textbf{Recommend}]
\begin{itemize}\scriptsize
  \item ``Austin: 18\% lower cost-per-hire, 2$\times$ the engineering talent pool, and 60\% of our top 20 clients within 2 hours.''
\end{itemize}
\end{shadedbox}
\end{column}
\end{columns}
\end{frame}

\begin{frame}{Why OAVR Matters for AI Prompting}
The quality of AI output depends on the quality of your \alert{question}, not the sophistication of your prompt.
\vspace{0.3cm}
\begin{columns}[T]
\begin{column}{0.45\textwidth}
\begin{shadedbox}[title=\textbf{Without OAVR}]
\begin{itemize}\small
  \item ``Analyze this data''
  \item AI produces generic summary statistics
  \item You get charts nobody asked for
\end{itemize}
\end{shadedbox}
\end{column}
\begin{column}{0.45\textwidth}
\begin{shadedbox}[title=\textbf{With OAVR}]
\begin{itemize}\small
  \item ``We're deciding between Austin and Denver. Compare cost-per-hire, talent pool, and client proximity.''
  \item AI produces exactly the analysis you need
\end{itemize}
\end{shadedbox}
\end{column}
\end{columns}
\vspace{0.3cm}
\alert{Orient first, then prompt.}  The OAVR framework ensures every analysis serves a decision.
\end{frame}

% ========================================
% SECTION 2: AI AS A THINKING PARTNER
% ========================================
\section{AI as a Thinking Partner}

\begin{frame}{AI as a Thinking Partner}
The hardest business decisions aren't data problems---they're \alert{thinking} problems.  AI isn't just an analyst.  It's a colleague who will argue with you.
\vspace{0.3cm}
\begin{itemize}
  \item \textbf{No agenda.}  Unlike colleagues, AI isn't protecting turf, managing up, or afraid to disagree with the boss.
  \item \textbf{Tireless devil's advocate.}  You can say ``argue against my preferred option''---no hurt feelings, no relationship cost.
  \item \textbf{Forces articulation.}  Explaining your reasoning to AI surfaces assumptions you didn't know you were making.
\end{itemize}
\vspace{0.3cm}
\alert{The conversation is the deliverable.}  You leave with a sharper view of the decision, not a report.
\end{frame}

\begin{frame}{Five Conversation Moves}
Specific things you can say to AI to sharpen your thinking.
\vspace{0.3cm}
\begin{center}\small
\begin{tabular}{@{}ll@{}}
\toprule
\textbf{Move} & \textbf{What you say} \\
\midrule
Expand options & ``What options am I not considering?'' \\
Challenge & ``Argue against my preferred choice.'' \\
Surface assumptions & ``What would need to be true for Option B to win?'' \\
Pre-mortem & ``If this fails, what's the most likely reason?'' \\
Reframe & ``How would a CFO / customer / regulator see this?'' \\
\bottomrule
\end{tabular}
\end{center}
\vspace{0.3cm}
These moves work because AI has no ego.  Use them early and often.
\end{frame}

\begin{frame}[shrink=15]{Example: Should We Cut Our Price?}
\begin{shadedbox}[title=\textbf{The Scenario}]\scriptsize
You're a VP at a SaaS company.  Your product is \$99/month.  A competitor just launched at \$79.  Your CEO wants a recommendation by Friday.
\end{shadedbox}
\vspace{0.1cm}
\small You open a conversation: \textit{``Help me think through whether we should drop our price to match.''}
\vspace{0.1cm}
\begin{columns}[T]
\begin{column}{0.5\textwidth}
\begin{shadedbox}[title=\textbf{AI expands your options}]
\begin{itemize}\scriptsize
  \item New \$79 tier with fewer features, keep \$99
  \item Add value at \$99 instead of cutting
  \item Hold price, offer selective retention discounts
\end{itemize}
\end{shadedbox}
\end{column}
\begin{column}{0.5\textwidth}
\begin{shadedbox}[title=\textbf{AI challenges assumptions}]
\begin{itemize}\scriptsize
  \item 20\% price cut needs 25\% more customers to break even
  \item Signals your product wasn't worth \$99
  \item Competitors can cut again---a race you can't win on price
\end{itemize}
\end{shadedbox}
\end{column}
\end{columns}
\vspace{0.1cm}
\centering\small Five minutes.  Three options instead of two, risks of each, and a failure scenario to monitor.  \alert{No data was harmed in the making of this decision.}
\end{frame}

% ========================================
% SECTION 3: STEELMANNING AND DEEPER THINKING
% ========================================
\section{Deeper Strategic Thinking}

\begin{frame}{Steelmanning}
\textbf{Steelmanning} = constructing the \alert{strongest possible version} of the opposing argument.  The opposite of a straw man.
\vspace{0.3cm}
\begin{columns}[T]
\begin{column}{0.5\textwidth}
\begin{shadedbox}[title=\textbf{How to Use It}]
\begin{itemize}\small
  \item ``I want to acquire Company X.  Give me the strongest possible argument \textit{against} this acquisition.''
  \item ``Steelman the bear case for our stock.''
  \item ``What's the best argument for the strategy I just rejected?''
\end{itemize}
\end{shadedbox}
\end{column}
\begin{column}{0.5\textwidth}
\begin{shadedbox}[title=\textbf{Why It Works}]
\begin{itemize}\small
  \item Forces you to address the \textit{real} counterargument, not a weak version
  \item Prevents confirmation bias
  \item If you can't defeat the steelman, you shouldn't proceed
\end{itemize}
\end{shadedbox}
\end{column}
\end{columns}
\end{frame}

\begin{frame}{Second-Order Thinking}
First-order thinking asks: ``What happens next?''  \alert{Second-order thinking} asks: ``And then what?''
\vspace{0.3cm}
\begin{columns}[T]
\begin{column}{0.5\textwidth}
\textbf{First-Order}
\begin{itemize}\small
  \item We cut price $\rightarrow$ we gain market share
  \item We lay off 10\% $\rightarrow$ costs go down
  \item We raise rates $\rightarrow$ inflation falls
\end{itemize}
\end{column}
\begin{column}{0.5\textwidth}
\textbf{Second-Order}
\begin{itemize}\small
  \item \ldots but margins shrink, competitors match, and we're stuck at \$79
  \item \ldots but morale drops, top talent leaves, and productivity falls
  \item \ldots but housing cools, banks tighten, and unemployment rises
\end{itemize}
\end{column}
\end{columns}
\vspace{0.3cm}
\textbf{Prompt:} ``Walk me through the second- and third-order consequences of [decision].  What happens after the obvious first effect?''
\end{frame}

\begin{frame}{Scenario Planning}
\textbf{Scenario planning} = building multiple plausible futures, not predicting the one ``right'' outcome.
\vspace{0.3cm}
\begin{columns}[T]
\begin{column}{0.5\textwidth}
\begin{shadedbox}[title=\textbf{The Approach}]
\begin{enumerate}\small
  \item Identify two key uncertainties (e.g., demand growth and regulatory risk)
  \item Build a 2$\times$2 matrix of scenarios
  \item Assess each scenario's impact and identify actions that perform well \textit{across} scenarios
\end{enumerate}
\end{shadedbox}
\end{column}
\begin{column}{0.5\textwidth}
\begin{shadedbox}[title=\textbf{AI's Role}]
\begin{itemize}\small
  \item ``Generate four scenarios for our business based on [uncertainty 1] and [uncertainty 2]''
  \item ``For each scenario, estimate the impact on our three-year revenue forecast''
  \item ``Which strategic option is most robust across all four scenarios?''
\end{itemize}
\end{shadedbox}
\end{column}
\end{columns}
\vspace{0.3cm}
Scenario planning + AI = stress-testing your strategy in minutes, not months.
\end{frame}

% ========================================
% SECTION 4: EFFECTIVE PROMPTING
% ========================================
\section{Effective Prompting}

\begin{frame}{Start with the Decision, Not the Data}
\begin{center}
{\Large\bfseries AI is most useful when you know what you're trying to decide.}
\end{center}

\vspace{0.5cm}

\begin{columns}[T]
\begin{column}{0.45\textwidth}
\begin{shadedbox}[title=\textbf{Weak}]
\begin{itemize}\small
  \item ``Show me the data''
  \item ``Make a chart of sales''
  \item ``Summarize this spreadsheet''
\end{itemize}
\end{shadedbox}
\end{column}
\begin{column}{0.45\textwidth}
\begin{shadedbox}[title=\textbf{Strong}]
\begin{itemize}\small
  \item ``Which regions should get more marketing budget?''
  \item ``Which product lines should we cut?''
  \item ``Is the discount strategy destroying margin?''
\end{itemize}
\end{shadedbox}
\end{column}
\end{columns}

\vspace{0.3cm}
Frame questions around the decision.  Follow-up questions cost nothing---just type.
\end{frame}

\begin{frame}{Prompting = Collaboration, Not Commands}
\begin{itemize}
  \item AI is not a search engine---it's a \alert{collaborator}; think of it as a capable colleague
  \item Do not try to craft the ``perfect prompt''---instead, have a \alert{conversation}
  \item The goal is \textit{iterative refinement}, not one-shot perfection
\end{itemize}

\vspace{0.3cm}

\begin{columns}[T]
\begin{column}{0.45\textwidth}
\begin{shadedbox}[title=\textbf{Don't Do This}]
\begin{itemize}\small
  \item Try to anticipate every edge case
  \item Give up when the first response is wrong
  \item Treat AI as a one-shot tool
\end{itemize}
\end{shadedbox}
\end{column}
\begin{column}{0.45\textwidth}
\begin{shadedbox}[title=\textbf{Do This Instead}]
\begin{itemize}\small
  \item Start with a rough request
  \item Refine based on the response
  \item Iterate until you get what you need
\end{itemize}
\end{shadedbox}
\end{column}
\end{columns}

\vspace{0.3cm}
\textbf{Ask the AI what it needs:} ``What information do you need from me to do this?''
\end{frame}

\begin{frame}{The Plan-Execute-Evaluate Cycle}
\begin{center}
\begin{tikzpicture}[scale=1.2]
  % Nodes
  \node[draw, rounded corners, fill=accentblue!20, minimum width=2.5cm, minimum height=1cm] (plan) at (0,2) {\textbf{Plan}};
  \node[draw, rounded corners, fill=accentblue!20, minimum width=2.5cm, minimum height=1cm] (execute) at (4,2) {\textbf{Execute}};
  \node[draw, rounded corners, fill=accentblue!20, minimum width=2.5cm, minimum height=1cm] (evaluate) at (2,0) {\textbf{Evaluate}};

  % Arrows
  \draw[->, thick, accentblue] (plan) -- (execute);
  \draw[->, thick, accentblue] (execute) -- (evaluate);
  \draw[->, thick, accentblue] (evaluate) -- (plan);
\end{tikzpicture}
\end{center}

\begin{itemize}
  \item \textbf{Plan:} Define the task, gather requirements, outline approach
  \item \textbf{Execute:} Let AI do the work with your guidance
  \item \textbf{Evaluate:} Check results, identify gaps, refine---then repeat
\end{itemize}
\end{frame}

\begin{frame}{Cross-Evaluation}
\begin{itemize}
  \item Start a \alert{new conversation} to evaluate the plan and output---even with the same model, a fresh context can catch errors
  \item Ask the new conversation to review, critique, or verify; different phrasing may reveal blind spots
  \item Consider using a \textit{different model} for additional perspective
\end{itemize}

\vspace{0.3cm}
\textbf{Example:} \textit{``I received this analysis from another session. Can you review it for errors or questionable assumptions?''}

\vspace{0.3cm}
This technique becomes critical in Module 8 (Verification and Governance).
\end{frame}

\end{document}
